%%%%%%%%%%%%
%
% $Autor: Chantal Crety $
% $Datum: 11.07.2025 $
% $Pfad: DemonstratorSchrittmotor\DemontageAnleitung\Einleitung.tex
% $Beschreibung: Einleitung DemontageAnleitung
%
%%%%%%%%%%%%

\chapter{Einleitung}

Herzlichen Glückwunsch zu Ihrem Schrittmotor Demonstrator! Diese Anleitung hilft Ihnen dabei, Ihren Demonstrator sicher und einfach zu demontieren. Die Schritt-für-Schritt-Anleitung zeigt Ihnen, wie Sie alle Teile fachgerecht auseinanderbauen können, ohne sie zu beschädigen. So können Sie Wartungsarbeiten oder Reparaturen selbstständig durchführen.


\vspace{0.5cm}
\textbf{Allgemein:}
	Die Baugruppe wird vormontiert geliefert, das heißt alle
	elektronischen Komponenten, inklusive des Sensors, befinden sich im montierten
	Zustand im Gehäuse des CNC-Tischs.  Die Demontage setzt keine besonderen Fachkenntnisse voraus – mit etwas handwerklichem Geschick und der Beachtung der Sicherheitshinweise gelingt Ihnen der Abbau problemlos.


	\vspace{2cm}
\fbox{\parbox{\textwidth}{
		\textbf{Hinweis:}\\
		Bitte bereiten Sie vor Beginn der Arbeiten Ihren Arbeitsplatz sorgfältig vor. Achten Sie während der Demontage auf eine schonende Handhabung aller elektronischen Bauteile, um Beschädigungen zu vermeiden.
		}}
	
	
	
	